\section{统一模型检验}

% 只有跨多个问题的统一检验才单独成章；单问检验应写回对应问题。
% 下列检验按题型和风险选择，不机械全部保留，也不默认五折交叉验证适用于所有数据。
\subsection{检验设计与数据隔离}

% 本节问题：要检验哪条核心结论，采用什么判据，哪些数据绝不参与选择？
% 结果：明确对照、指标、阈值和训练/验证/测试或时空分组边界。
说明要核验的结论、指标、对照、判据和数据隔离方式。IID、分组、时序和空间数据分别采用匹配的划分；模型选择只看训练/验证信息，最终测试结果不得反复参与调参。

\subsection{误差、残差与独立互验}

% 本节问题：误差是否有结构性偏差，核心答案能否被独立方法或精确/理论边界复核？
% 结果：给出误差分布、配对差、区间/离散程度、失败案例或可行性结论。
根据任务报告误差/残差、理论或精确解互验、经典基线比较、约束可行率或收敛性。随机结果给重复次数、种子口径、均值与离散程度。

\subsection{敏感性、不确定性与情景}

% 本节问题：观测、参数和情景变化分别怎样影响结论，何时需要换策略？
% 结果：报告有依据的扰动范围、区间、排序/方案是否改变及触发条件。
区分观测误差、参数估计误差和未来情景变化。只有输入分布或情景有依据时才做蒙特卡洛或情景矩阵；同时报告区间、方案是否改变以及触发换策略的条件。

\subsection{检验结论与限制}

% 本节问题：哪些主张被支持，哪些仍只能称风险或待验证，结论边界在哪里？
% 接口：把通过/未通过的结论逐项回写摘要、正文、评价与结论章。
报告可复算结果、失败场景和结论边界，不写没有实验支撑的“稳定”“显著”或泛化结论。
